\chapter{Dimensionamiento}\label{cap:dim}

\section{Datos de diseño}
\begin{table}[H]
\centering\small
\caption{Datos de partida.}\label{tab:datos}
\begin{tabular}{@{}llr@{}}
\toprule
\textbf{Magnitud} & \textbf{Símbolo} & \textbf{Valor}\\\midrule
Masa total (CanSat + contenedor)       & $m$        & \SI{0.500}{kg}\\
Peso                                   & $W$        & \SI{4.905}{N}\\
Densidad del aire                      & $\rho$     & \SI{1.225}{kg/m^3}\\
Viscosidad cinemática                  & $\nu$      & \SI{1.5e-5}{m^2/s}\\
Radio del rotor                        & $R$        & \SI{0.200}{m}\\
Radio de la bisagra                    & $e$        & \SI{0.044}{m}\\
Largo de pala                          & $L_b=R-e$  & \SI{0.156}{m}\\
Número de palas                        & $B$        & 4\\
Cuerda                                 & $c$        & \SI{0.045}{m}\\
Área del disco                         & $A_R=\pi R^2$ & \SI{0.1257}{m^2}\\
Solidez                                & $\sigma=Bc/\pi R$ & 0,286\\
Diámetro del drogue                    & $D_d$      & \SI{0.25}{m}\\
Área del drogue                        & $A_d$      & \SI{0.0491}{m^2}\\
Coeficiente del drogue                 & $C_{D,d}$  & 0,80\\
Área frontal del cuerpo (Ø\,95)        & $A_c$      & \SI{0.00709}{m^2}\\
Coeficiente del cuerpo                 & $C_{D,c}$  & 0,80\\
\bottomrule
\end{tabular}
\end{table}

%=====================================================================
\section{Etapa 1: drogue}\label{sec:dim-drogue}
Con la ecuación~\eqref{eq:equilibrio} sin rotor (las palas van plegadas) y la masa
completa:

\begin{figure}[H]
\centering
\begin{tikzpicture}
\begin{axis}[width=0.9\textwidth, height=6.6cm, xmin=15, xmax=40, ymin=6, ymax=24,
  xlabel={Diámetro del drogue $D_d$ [cm]}, ylabel={Velocidad etapa 1 [m/s]},
  legend style={font=\footnotesize, at={(0.98,0.98)}, anchor=north east}]
  \fill[ccuerda!12] (axis cs:15,12) rectangle (axis cs:40,18);
  \node[font=\footnotesize, ccuerda!70!black, anchor=south east] at (axis cs:39.5,12.2) {\Req{C3}: \SIrange{12}{18}{m/s}};
  \addplot[cfijo, thick, dashed, domain=15:40, samples=60] {sqrt(4.905/(0.6125*(0.75*pi*(x/100)^2/4+0.8*0.007088)))};
  \addlegendentry{$C_{D,d}=0{,}75$}
  \addplot[cfijo, very thick, domain=15:40, samples=60] {sqrt(4.905/(0.6125*(0.80*pi*(x/100)^2/4+0.8*0.007088)))};
  \addlegendentry{$C_{D,d}=0{,}80$}
  \addplot[cfijo, thick, dotted, domain=15:40, samples=60] {sqrt(4.905/(0.6125*(0.90*pi*(x/100)^2/4+0.8*0.007088)))};
  \addlegendentry{$C_{D,d}=0{,}90$}
  \addplot[only marks, mark=*, mark size=3pt, cfreno] coordinates {(25,13.35)};
  \node[cfreno, font=\footnotesize, anchor=south west] at (axis cs:25.4,13.5) {Ø\,25\,cm: \SI{13.4}{m/s}};
\end{axis}
\end{tikzpicture}
\caption{Velocidad de la primera etapa según el diámetro del drogue (incluye el cuerpo).}
\label{fig:drogue}
\end{figure}

\begin{table}[H]
\centering\small
\caption{Velocidad de la etapa 1 (drogue + cuerpo, \SI{500}{g}).}\label{tab:drogue}
\begin{tabular}{@{}cccccc@{}}
\toprule
$D_d$ [cm] & $A_d$ [m$^2$] & $C_{D,d}=0{,}75$ & $C_{D,d}=0{,}80$ & $C_{D,d}=0{,}90$ & \Req{C3}\\\midrule
24 & 0,0452 & 14,2 & 13,8 & 13,1 & Cumple\\
\textbf{25} & \textbf{0,0491} & \textbf{13,7} & \textbf{13,4} & \textbf{12,7} & \textbf{Cumple}\\
27 & 0,0573 & 12,8 & 12,5 & 11,8 & No cumple con $C_D$ alto\\
30 & 0,0707 & 11,7 & 11,3 & 10,8 & No cumple\\
\bottomrule
\end{tabular}
\end{table}

\begin{resultado}[Drogue elegido]
Ø\,\SI{25}{cm}, cruz o circular plano, $C_D\approx0{,}8$: $\Vd\approx\SI{13.4}{m/s}$.
Queda dentro de \Req{C3} para todo el rango de $C_D$ esperable. Un drogue mayor mejoraría
la etapa 2, pero a Ø\,\SI{27}{cm} la etapa 1 ya puede caer por debajo de \SI{12}{m/s}.
\end{resultado}

\paragraph{Carga de apertura.} Con la ecuación~\eqref{eq:apertura}, $C_D S=\SI{0.0393}{m^2}$:
a $V_{\text{ap}}=\SI{15}{m/s}$, $F_o\approx\SI{8}{N}$; a \SI{25}{m/s}, \SI{23}{N}; a
\SI{35}{m/s}, \SI{44}{N}. El CanSat sale del cohete cerca del apogeo, a baja velocidad, así
que \SI{25}{m/s} es un caso conservador. Se especifica una cuerda de Dyneema de
\SI{1.5}{mm} (carga de rotura del orden de \SI{1}{kN}) y un destorcedor de
$\ge\SI{200}{N}$: factor de seguridad mayor que 4 aun a \SI{35}{m/s}.

%=====================================================================
\section{Etapa 2: rotor}\label{sec:dim-rotor}

\subsection{Velocidad de descenso}
Con la ecuación~\eqref{eq:equilibrio}, cada término aporta $k=\tfrac12\rho\,C\,A$ en
\si{N/(m/s)^2}: rotor $k_R=0{,}0770\,\CR$, drogue $k_d=0{,}0241$, cuerpo $k_c=0{,}0035$.

\begin{table}[H]
\centering\small
\caption{Velocidad estabilizada de la etapa 2 (\SI{500}{g}, $R=\SI{0.20}{m}$, drogue Ø\,25).}\label{tab:rotor}
\begin{tabular}{@{}lccc@{}}
\toprule
\textbf{Caso} & $\boldsymbol{\CR}$ & \textbf{Con drogue} & \textbf{Drogue sin aporte}\\\midrule
Rotor pobre                    & 1,0 & 6,85 & 7,81\\
Rotor real típico              & 1,1 & 6,61 & 7,46\\
\textbf{Diseño}                & \textbf{1,2} & \textbf{6,40} & \textbf{7,15}\\
Rotor muy bueno                & 1,3 & 6,20 & 6,88\\
\bottomrule
\multicolumn{4}{@{}l}{\footnotesize Velocidades en m/s. \Req{C4} admite hasta \SI{8}{m/s}.}
\end{tabular}
\end{table}

\begin{figure}[H]
\centering
\begin{tikzpicture}
\begin{axis}[width=0.92\textwidth, height=7.2cm, xmin=0.10, xmax=0.32, ymin=2, ymax=14,
  xlabel={Radio del rotor $R$ [m]}, ylabel={Velocidad etapa 2 [m/s]},
  xticklabel style={/pgf/number format/fixed, /pgf/number format/precision=2},
  legend columns=2, legend style={font=\footnotesize, at={(0.5,-0.22)}, anchor=north, cells={anchor=west}}]
  \fill[ccuerda!12] (axis cs:0.10,2) rectangle (axis cs:0.32,8);
  \node[font=\footnotesize, ccuerda!70!black, anchor=south west] at (axis cs:0.255,2.2) {\Req{C4}: \SIrange{2}{8}{m/s}};
  \fill[cfreno!8] (axis cs:0.20,2) rectangle (axis cs:0.32,14);
  \node[font=\footnotesize, cfreno, anchor=north west, align=left] at (axis cs:0.203,13.7) {Fuera de la\\ restricción $R\le0{,}20$};
  \addplot[crot, very thick, domain=0.10:0.32, samples=60] {sqrt(4.905/(0.6125*1.2*pi*x^2+0.02405+0.003473))};
  \addlegendentry{$\CR=1{,}2$, con drogue}
  \addplot[crot, thick, dashed, domain=0.10:0.32, samples=60] {sqrt(4.905/(0.6125*1.0*pi*x^2+0.02405+0.003473))};
  \addlegendentry{$\CR=1{,}0$, con drogue}
  \addplot[cfijo, thick, domain=0.10:0.32, samples=60] {sqrt(4.905/(0.6125*1.2*pi*x^2+0.003473))};
  \addlegendentry{$\CR=1{,}2$, drogue sin aporte}
  \addplot[cfijo, thick, dashed, domain=0.10:0.32, samples=60] {sqrt(4.905/(0.6125*1.0*pi*x^2+0.003473))};
  \addlegendentry{$\CR=1{,}0$, drogue sin aporte (peor caso)}
  \draw[black!60, dashed] (axis cs:0.10,5) -- (axis cs:0.32,5);
  \node[font=\scriptsize, anchor=south west] at (axis cs:0.10,5) {objetivo \SI{5}{m/s}};
  \addplot[only marks, mark=*, mark size=3pt, cfreno] coordinates {(0.20,6.40)};
  \node[cfreno, font=\footnotesize, anchor=north east, fill=white, inner sep=1pt] at (axis cs:0.197,6.1) {diseño: \SI{6.4}{m/s}};
  \addplot[only marks, mark=square*, mark size=2.5pt, cgris] coordinates {(0.20,7.81)};
  \node[cgris, font=\footnotesize, anchor=west, fill=white, inner sep=1pt] at (axis cs:0.205,7.3) {peor caso: \SI{7.8}{m/s}};
\end{axis}
\end{tikzpicture}
\caption{Velocidad de la segunda etapa en función del radio. Con $R=\SI{0.20}{m}$ todos los
casos quedan bajo \SI{8}{m/s}; llegar a \SI{5}{m/s} exactos pediría $R\approx\SI{0.27}{m}$.}
\label{fig:VR}
\end{figure}

\subsection{Verificación con la curva empírica}
Con $\Vr=\SI{6.40}{m/s}$, el empuje del rotor es $T=W-(k_d+k_c)\Vr^2=\SI{3.78}{N}$ y la
carga de disco $T/A_R=\SI{30.1}{N/m^2}$. Por la ecuación~\eqref{eq:vh},
$\vh=\SI{3.50}{m/s}$ y $\Vr/\vh=1{,}83$. El polinomio~\eqref{eq:poly} da
$\vi=\SI{6.34}{m/s}$, es decir un flujo neto a través del disco de
$\SI{-0.06}{m/s}\approx0$: el punto de diseño es prácticamente la autorrotación ideal, lo
que confirma que $\CR=1{,}2$ es consistente con los datos empíricos
(figura~\ref{fig:estados}).

\subsection{Contraste con elemento de pala}
Un modelo de elemento de pala con iteración de velocidad inducida (placa plana, 4 palas de
\SI{35}{mm}, paso $\ang{-2}$), resuelto para par nulo, dio \SI{6.7}{m/s} con drogue y
\SI{7.6}{m/s} sin él, a \SIrange{2300}{2800}{rpm}. Coincide con la tabla~\ref{tab:rotor}
dentro del 5\,\%. Para perfiles muy curvados ese modelo pierde validez (predice $\CR>1{,}4$,
físicamente imposible), por lo que para la placa del 7\,\% se usa el límite
$\CR\approx1{,}2$ y la ecuación~\eqref{eq:omega} para las rpm.

%=====================================================================
\section{Perfil y velocidad de giro}\label{sec:dim-perfil}

\begin{table}[H]
\centering\small
\caption{Velocidad de giro según el perfil, con $T=\SI{3.78}{N}$ (ecuación~\ref{eq:omega}).}\label{tab:rpm}
\begin{tabular}{@{}lcccccc@{}}
\toprule
\textbf{Perfil} & $\boldsymbol{c}$ [mm] & $\boldsymbol{\bar C_l}$ & $\boldsymbol{\Omega}$ [rad/s] & \textbf{rpm} & $\boldsymbol{\Omega R}$ [m/s] & $\boldsymbol{Re}$ \textbf{punta}\\\midrule
Placa plana 3\,\%            & 35 & 0,3 & 236 & 2253 & 47,2 & \num{110000}\\
Fondo plano                  & 35 & 0,6 & 167 & 1593 & 33,4 & \num{78000}\\
Placa curvada 7\,\%          & 35 & 0,9 & 136 & 1301 & 27,2 & \num{64000}\\
\textbf{Placa curvada 7\,\%} & \textbf{45} & \textbf{0,9} & \textbf{120} & \textbf{1147} & \textbf{24,0} & \textbf{\num{72000}}\\
Placa curvada 7\,\%          & 50 & 1,0 & 108 & 1033 & 21,6 & \num{72000}\\
\bottomrule
\end{tabular}
\end{table}

\begin{figure}[H]
\centering
\begin{tikzpicture}
\begin{axis}[width=0.88\textwidth, height=6.4cm, xmin=0.008, xmax=0.052, ymin=500, ymax=3000,
  xlabel={Producto $c\,\bar C_l$ [m]}, ylabel={Velocidad de giro [rpm]},
  xticklabel style={/pgf/number format/fixed, /pgf/number format/precision=3},
  scaled x ticks=false]
  \addplot[crot, very thick, domain=0.008:0.052, samples=80] {9.5493*sqrt(584.6/x)};
  \addplot[only marks, mark=*, mark size=2.5pt, cfijo] coordinates
    {(0.0105,2253) (0.021,1593) (0.0315,1301) (0.0405,1147) (0.050,1033)};
  \node[font=\scriptsize, anchor=south west] at (axis cs:0.0105,2253) {placa plana};
  \node[font=\scriptsize, anchor=south west] at (axis cs:0.021,1593) {fondo plano};
  \node[font=\scriptsize, anchor=south west] at (axis cs:0.0315,1301) {curvada, $c=35$};
  \node[font=\scriptsize, anchor=south west, cfreno] at (axis cs:0.0405,1147) {\textbf{elegida}};
  \node[font=\scriptsize, anchor=north] at (axis cs:0.050,1000) {$c=50$};
\end{axis}
\end{tikzpicture}
\caption{La velocidad de giro baja con la raíz de $c\,\bar C_l$ (con el mismo empuje y la
misma velocidad de caída).}
\label{fig:rpm}
\end{figure}

\begin{resultado}[Perfil elegido]
Placa curvada al 7\,\%, cuerda \SI{45}{mm}, $\bar C_l\approx0{,}9$:
$\Omega\approx\SI{120}{rad/s}$ ($\approx\SI{1150}{rpm}$), velocidad de punta
\SI{24}{m/s}, $Re\approx\num{36000}$ en media pala y \num{72000} en la punta. La velocidad
de caída no cambia respecto de la placa plana; lo que cambia es la rpm, a la mitad. Las
rpm tienen una incertidumbre de $\pm$20\,\% y se miden en el ensayo con ventilador.
\end{resultado}

%=====================================================================
\section{Geometría y fabricación de la pala}\label{sec:pala}

\paragraph{Curvatura y molde.} Un arco circular de cuerda $c$ y flecha $s$ tiene radio
\begin{equation}
  R_c = \frac{c^2/4 + s^2}{2s}, \qquad s = R_c - \sqrt{R_c^2 - c^2/4}. \label{eq:flecha}
\end{equation}
Para $c=\SI{45}{mm}$ y 7\,\% ($s=\SI{3.15}{mm}$), $R_c=\SI{81.9}{mm}$. Un caño de PVC de
Ø\,\SI{160}{mm} ($R_c=\SI{80}{mm}$) da $s=\SI{3.23}{mm}$, o sea 7,2\,\%: es el molde.

\paragraph{Material y masa.} Placa de carbono de \SI{0.5}{mm} ($\rho\approx\SI{1550}{kg/m^3}$):
$156\times45\times0{,}5$\,mm da \SI{5.4}{g}; con portapala y refuerzo de raíz,
$m_b\approx\SI{7.4}{g}$. Alternativa: balsa de \SIrange{1}{1.5}{mm} laminada con tejido de
carbono.

\paragraph{Detalles.} Borde de ataque redondeado; borde de fuga afilado; centro de masa de
la sección cerca del 25\,\% de la cuerda (si queda más atrás, agregar una varilla de
carbono de \SI{1}{mm} en el borde de ataque) para evitar flameo; sin torsión en la primera
versión.

\paragraph{Paso.} Se fija en el portapala. Por la sección~\ref{sec:arranque} debe ser
$\theta\le\ang{0}$; se fabrican portapalas de $\ang{-6}$, $\ang{-4}$, $\ang{-2}$ y $\ang{0}$
y se elige en el ensayo con ventilador.

%=====================================================================
\section{Plegado dentro de Ø\,95}\label{sec:plegado}
Las palas cuelgan plegadas contra el cuerpo con la cuerda en dirección tangencial y la
cara cóncava hacia adentro. Al girar \ang{90} hacia afuera sobre la bisagra, esa cara
pasa a ser la inferior: la curvatura queda en el sentido correcto para sustentar.

Para una pala de cuerda $c$ y flecha $s$ cuyo punto medio exterior está a radio $r_m$, los
bordes quedan a
\begin{equation}
  r_{\text{borde}} = \sqrt{(r_m-s)^2 + (c/2)^2} \le \SI{47.5}{mm}
  \;\Rightarrow\; r_m \le s + \sqrt{47{,}5^2 - 22{,}5^2} = \SI{44.98}{mm}.
\end{equation}
Con espesor de \SI{0.5}{mm}, la cara interior queda a \SI{44.5}{mm} del eje: el cuerpo
puede tener \textbf{Ø\,\SI{88}{mm}} (radio \SI{44}{mm}) con \SI{0.5}{mm} de juego. Cada
pala ocupa $2\arcsin(22{,}5/47{,}5)=\ang{56.5}$ del perímetro; las cuatro, \ang{226}.
Quedan cuatro huecos de \ang{33.5} ($\approx\SI{27}{mm}$ de arco) para
los orificios de los interruptores y los paneles solares.

%=====================================================================
\section{Batimiento y conicidad}\label{sec:dim-batimiento}
Con $L=T/B=\SI{0.945}{N}$ por pala aplicada a $0{,}75R$, el brazo respecto de la bisagra es
$0{,}75R-e=\SI{0.106}{m}$ y $M_L=\SI{0.100}{N.m}$; el peso de la pala aporta
$M_W=\SI{0.006}{N.m}$. Con $m_b=\SI{7.4}{g}$ y $\Omega=\SI{120}{rad/s}$, la rigidez
centrífuga es $m_b\Omega^2\overline{r(r-e)}=\SI{1.24}{N.m/rad}$ y, por~\eqref{eq:conicidad},
\begin{equation*}
  \beta_0 \approx \frac{0{,}100-0{,}006}{1{,}24} = \SI{0.076}{rad} \approx \ang{4.4}.
\end{equation*}
El número de Lock~\eqref{eq:lock} es $\gamma\approx8{,}3$ (dentro del rango de los
helicópteros) y la frecuencia de batimiento~\eqref{eq:nubeta},
$\nu_\beta/\Omega=1{,}19$: no hay resonancia con la excitación de 1 vez por vuelta.

\begin{resultado}[Topes de la bisagra]
Tope superior a \ang{+10}: en vuelo la pala trabaja libre en $\approx\ang{4.4}$ y el tope
solo actúa en el arranque y en ráfagas. Sin tope inferior en el rango de vuelo (la pala
debe poder plegarse \ang{90} hacia abajo para el almacenamiento). El resorte de apertura
($\approx\SI{0.005}{N.m}$) desplaza el equilibrio menos de \ang{0.3}.
\end{resultado}

%=====================================================================
\section{Arranque}\label{sec:arranque}
En la liberación el CanSat cae a $\Vd\approx\SI{13.4}{m/s}$. Con la
ecuación~\eqref{eq:Q0}:

\begin{table}[H]
\centering\small
\caption{Par de arranque con el rotor quieto ($V=\SI{13.4}{m/s}$).}\label{tab:arranque}
\begin{tabular}{@{}lccccc@{}}
\toprule
Paso $\theta$ & $\ang{-6}$ & $\ang{-4}$ & $\ang{-2}$ & $\ang{0}$ & $\ang{+2}$\\\midrule
$Q_0$ [N\,m] & 0,078 & 0,052 & 0,026 & 0 & $-0{,}026$\\
Arranca solo & Sí & Sí & Sí & No & No (gira al revés)\\
\bottomrule
\end{tabular}
\end{table}

La inercia polar del rotor es $I\approx\SI{5.1e-4}{kg.m^2}$ y su energía a régimen,
$\tfrac12 I\Omega^2=\SI{3.7}{J}$. El par disminuye a medida que el rotor acelera (se anula
en el equilibrio); con un par medio de \SIrange{0.01}{0.02}{N.m}, el tiempo de arranque es
$t=I\Omega/\bar Q\approx\SIrange{3}{6}{s}$, durante los cuales se pierden unos
\SIrange{30}{60}{m} de altura. Con apogeo de \SI{700}{m}, la liberación ocurre a
\SI{560}{m}: el margen es amplio.

%=====================================================================
\section{Cargas y resistencia de los componentes}\label{sec:cargas}

\begin{table}[H]
\centering\small
\caption{Cargas en el rotor y verificación.}\label{tab:cargas}
\begin{tabularx}{\textwidth}{@{}lcY@{}}
\toprule
\textbf{Carga} & \textbf{Valor} & \textbf{Verificación}\\\midrule
Fuerza centrífuga por pala ($m_b r_g\Omega^2$) & \SI{13.1}{N} & Pasador de acero Ø\,2 en doble corte: $\tau=\SI{2.1}{MPa}$, frente a $\approx\SI{300}{MPa}$\\
Sustentación total & \SI{3.78}{N} & Rodamientos 688ZZ: capacidad dinámica del orden de \SI{1}{kN}; vida ilimitada\\
Momento en la raíz (tope en el arranque) & $\le\SI{0.10}{N.m}$ & Carbono \SI{0.5}{mm}, $45\times0{,}5^2/6=\SI{1.9}{mm^3}$: \SI{53}{MPa} frente a $\approx\SI{500}{MPa}$; reforzar raíz con portapala de \SIrange{15}{20}{mm}\\
Desbalance \SI{0.1}{g} a \SI{0.12}{m} & \SI{0.17}{N} & Aceptable (3,5\,\% del peso)\\
Desbalance \SI{1}{g} a \SI{0.12}{m} & \SI{1.73}{N} & Inaceptable: balancear a $\pm\SI{0.1}{g}$\\
Inercia de una pala a 30\,G (\Req{S5}) & \SI{2.2}{N} & Pestaña de acero inoxidable de \SI{0.5}{mm} y diente del anillo; sin riesgo\\
Apertura del drogue (\SI{25}{m/s}) & \SI{23}{N} & Va al anclaje en la estructura, no a los rodamientos\\
\bottomrule
\end{tabularx}
\end{table}

%=====================================================================
\section{Estructura}\label{sec:estructura}
Tubo de PETG Ø\,88 con pared de \SI{2}{mm} (\Req{S6}). Bajo 30\,G con \SI{0.5}{kg}, la
fuerza axial es \SI{147}{N} y la tensión
$\sigma=\SI{0.25}{MPa}$. La tensión crítica de pandeo de un cilindro de pared delgada,
con factor de reducción 0,3 por imperfecciones~\cite{timoshenko1961},
\begin{equation}
  \sigma_{cr} = 0{,}3\,\frac{E\,t}{R_m\sqrt{3(1-\nu^2)}}
  = 0{,}3\,\frac{\SI{2}{GPa}\cdot\SI{2}{mm}}{\SI{46.5}{mm}\cdot1{,}60} \approx \SI{16}{MPa},
\end{equation}
da un margen de 64 veces. La pared no flexiona ni se deforma bajo las cargas del anexo; el
punto débil son las uniones (tapas y anillos), que se atornillan con insertos roscados.

%=====================================================================
\section{Tiempos de descenso}
\begin{table}[H]
\centering\small
\caption{Tiempos con liberación al 80\,\% del apogeo ($\Vd=\SI{13.4}{m/s}$,
$\Vr=\SI{6.4}{m/s}$; se desprecia el arranque).}\label{tab:tiempos}
\begin{tabular}{@{}cccccc@{}}
\toprule
\textbf{Apogeo} [m] & \textbf{Liberación} [m] & \textbf{Etapa 1} [s] & \textbf{Etapa 2} [s] & \textbf{Total} [s] & \textbf{Video necesario}\\\midrule
500  & 400 & 7,5  & 62,5  & 70  & $\approx$\,2 min\\
700  & 560 & 10,4 & 87,5  & 98  & $\approx$\,2 min\\
1000 & 800 & 15,0 & 125,0 & 140 & $\approx$\,3 min\\
\bottomrule
\end{tabular}
\end{table}

%=====================================================================
\section{Presupuestos}\label{sec:masa}

\subsection{Masa}
\begin{table}[H]
\centering\small
\caption{Presupuesto de masa estimado.}\label{tab:masa}
\begin{tabular}{@{}lrr@{}}
\toprule
\textbf{Componente} & \textbf{Config. A} [g] & \textbf{Config. B} [g]\\\midrule
Cuerpo PETG Ø\,88 (B: Ø\,83), pared \SI{2}{mm}, \SI{205}{mm} & 141 & 133\\
Tapas y anillos internos                                   & 25 & 25\\
Camisa (contenedor) PETG Ø\,95/91, \SI{205}{mm}            & --- & 152\\
Rotor: 4 palas con portapala                               & 30 & 30\\
Rotor: cubo, eje, rodamientos, tuerca, resortes, sensor    & 24 & 24\\
Drogue, copa, cuerda y destorcedor                         & 18 & 18\\
Traba: anillo interno, servo, pestañas                       & 16 & 16\\
Electrónica (MCU, IMU, barómetro, GPS, radio, SD, PCB)     & 47 & 47\\
Batería 18650 y portapila                                  & 55 & 55\\
Cámaras (2)                                                & 20 & 20\\
Baliza con pila propia                                     & 8 & 8\\
Paneles solares (2)                                        & 8 & 8\\
Cables, tornillos e insertos                               & 25 & 25\\\midrule
\textbf{Subtotal}                                          & \textbf{417} & \textbf{561}\\
Lastre (abajo, ajustable)                                  & 83 & ---\\
\textbf{Total frente a \SI{500 \pm 10}{g}}                 & \textbf{500} & \textbf{561 (excede 51)}\\
\bottomrule
\end{tabular}
\end{table}

La configuración A tiene \SI{83}{g} de margen, que se usa como lastre en el fondo (baja el
centro de masa y estabiliza el descenso). La B excede el límite; solo sería viable con una
camisa no impresa (cartón o fibra, a la que no aplica \Req{S6}) o con ventanas grandes.

\subsection{Largo}
\begin{table}[H]
\centering\small
\caption{Distribución del largo (configuración A, medido desde el fondo).}\label{tab:largo}
\begin{tabular}{@{}lcc@{}}
\toprule
\textbf{Tramo} & \textbf{Desde--hasta} [mm] & \textbf{Largo} [mm]\\\midrule
Tapa inferior y cámara nadir & 0--3 & 3\\
Baliza, lastre y cámara nadir & 3--40 & 37\\
Anillo de traba y servo (pestañas a \SI{60}{mm}) & 40--70 & 30\\
Batería 18650 vertical & 70--140 & 70\\
Electrónica (pila de PCB) & 140--195 & 55\\
Tapa superior y cámara cenit & 195--205 & 10\\
Cubo del rotor (bisagras a \SI{213}{mm}) & 205--222 & 17\\
Copa del drogue (contenedor mínimo) & 222--250 & 28\\\midrule
\textbf{Total} & & \textbf{250}\\
\bottomrule
\end{tabular}
\end{table}

\subsection{Energía}
\begin{table}[H]
\centering\small
\caption{Consumo estimado.}\label{tab:energia}
\begin{tabular}{@{}lr@{}}
\toprule
\textbf{Consumo} & \textbf{Potencia} [W]\\\midrule
Microcontrolador & 0,30\\
Radio (ráfagas a 1\,Hz, promedio) & 0,25\\
GPS & 0,10\\
Sensores (IMU, barómetro, corriente, hall) & 0,02\\
Cámaras (2) & 1,80\\
Servo en reposo & 0,03\\\midrule
\textbf{Total} & \textbf{2,50}\\
\bottomrule
\end{tabular}
\end{table}
Una celda 18650 de \SI{3000}{mAh} y \SI{3.6}{V} guarda \SI{10.8}{Wh}; con 85\,\% de
eficiencia de conversión da \SI{9.2}{Wh}, es decir unas \SI{3.7}{h} con todo encendido. Si
las cámaras se encienden recién en \texttt{ASCENT}, la espera en la plataforma puede superar
las \SI{10}{h}.
